\documentclass[reqno]{amsart}
\usepackage[T1]{fontenc}
\usepackage{amsmath,amssymb,amsthm}
\usepackage{xr-hyper}
\usepackage[hidelinks]{hyperref}
\externaldocument[hk-]{20261125T016500--hecke-operators-and-theta-lifts__workbook}
\externaldocument[tt-]{20261125T004100--tate-twists-and-sha__workbook}
\externaldocument[gs-]{20261126T020430--the-minus-part-by-gauss-sums__workbook}
\externaldocument[cg-]{20261125T014100--congruence-modules__workbook}
\externaldocument[ec-]{20261125T014500--eisenstein-congruences-and-selmer-classes__workbook}
\externaldocument[pc-]{20261125T014700--pseudocharacters-and-the-length-bound__workbook}
\externaldocument[gr-]{20261125T011000--galois-representations-black-box__workbook}
\externaldocument[ob-]{20261126T040100--spaces-forms-and-twists__workbook}
\externaldocument[dp-]{20261126T030900--denominators-of-projections-to-lifts__workbook}
\externaldocument[n2-]{20261126T060500--auxiliary-primes-even-classes-at-no-cost__workbook}

% Chapter number; \ref{chap} gives it.
\newcounter{chap}
% Chapter numbers, generated from 20261125T012000--grh-workbook__guide.org
% by scripts/make-chapters.py.  Do not edit: change the order of the
% index and rerun the script (./preview-site.sh and CI do this).
%
% Usage in a sheet, after \newcounter{chap}:
%   % Chapter numbers, generated from 20261125T012000--grh-workbook__guide.org
% by scripts/make-chapters.py.  Do not edit: change the order of the
% index and rerun the script (./preview-site.sh and CI do this).
%
% Usage in a sheet, after \newcounter{chap}:
%   % Chapter numbers, generated from 20261125T012000--grh-workbook__guide.org
% by scripts/make-chapters.py.  Do not edit: change the order of the
% index and rerun the script (./preview-site.sh and CI do this).
%
% Usage in a sheet, after \newcounter{chap}:
%   \input{chapters}\chapnumber{<sheet ID>}
\makeatletter
\providecommand\chapters@set[2]{\@namedef{chapters@#1}{#2}}
\chapters@set{20261125T001000}{1}% zeros-to-ramp-sums
\chapters@set{20261125T001100}{2}% inert-primes-and-estimate-b
\chapters@set{20261125T005000}{3}% galois-fourier-and-averaging
\chapters@set{20261125T005500}{4}% columns-from-a-sphere
\chapters@set{20261125T000000}{5}% slope-method
\chapters@set{20261126T030100}{6}% a-twist-for-each-prime
\chapters@set{20261125T003000}{7}% the-chain-and-the-square-root
\chapters@set{20261125T010000}{8}% rank-lemma-and-zero-estimate
\chapters@set{20261125T010500}{9}% poincare-series-and-near-orthogonality
\chapters@set{20261124T223000}{10}% elliptic-curves-for-the-proof
\chapters@set{20261125T009500}{11}% fock-model-and-theta-intertwiners
\chapters@set{20261124T223100}{12}% complex-multiplication-and-inert-reduction
\chapters@set{20261125T008000}{13}% one-cm-determinant
\chapters@set{20261124T214500}{14}% supersingular-theta-divisibility
\chapters@set{20261124T214600}{15}% the-isotropic-plane
\chapters@set{20261125T013700}{16}% pairings-on-a-formal-group-of-height-two
\chapters@set{20261125T013500}{17}% hecke-nondegeneracy-on-the-tree
\chapters@set{20261126T030200}{18}% the-bound-is-attained
\chapters@set{20261125T009000}{19}% kronecker-limit-and-chowla-selberg
\chapters@set{20261125T002000}{20}% definite-theta-lifts
\chapters@set{20261125T015500}{21}% siegel-weil-in-the-smallest-case
\chapters@set{20261125T002100}{22}% definite-theta-lifts-and-norms
\chapters@set{20261125T016500}{23}% hecke-operators-and-theta-lifts
\chapters@set{20261125T016700}{24}% divisibility-along-the-chain
\chapters@set{20261125T004000}{25}% denominators-and-sha
\chapters@set{20261125T004100}{26}% tate-twists-and-sha
\chapters@set{20261126T020430}{27}% the-minus-part-by-gauss-sums
\chapters@set{20261125T014100}{28}% congruence-modules
\chapters@set{20261125T014500}{29}% eisenstein-congruences-and-selmer-classes
\chapters@set{20261125T014700}{30}% pseudocharacters-and-the-length-bound
\chapters@set{20261125T011000}{31}% galois-representations-black-box
\chapters@set{20261126T030300}{32}% the-data
\chapters@set{20261126T020210}{33}% bessel-poincare-series-and-their-gram-operators
\chapters@set{20261126T030400}{34}% the-seesaw-identity
\chapters@set{20261126T030500}{35}% divisibility-of-the-composite-lift
\chapters@set{20261126T030600}{36}% theorem-a-short-forms
\chapters@set{20261126T030700}{37}% the-determinant-of-a-chain
\chapters@set{20261126T030800}{38}% the-cost-of-one-lift
\chapters@set{20261126T030900}{39}% denominators-of-projections-to-lifts
\chapters@set{20261126T031000}{40}% twists-with-few-congruences
\chapters@set{20261126T031100}{41}% theorem-c-and-the-main-inequality
\providecommand\chapnumber[1]{%
  \@ifundefined{chapters@#1}%
    {\PackageError{chapters}{Sheet #1 is not listed in the index}%
      {List it in 20261125T012000--grh-workbook__guide.org\MessageBreak
       and run scripts/make-chapters.py.}}%
    {\edef\thechap{\@nameuse{chapters@#1}}}}
\makeatother
\chapnumber{<sheet ID>}
\makeatletter
\providecommand\chapters@set[2]{\@namedef{chapters@#1}{#2}}
\chapters@set{20261125T001000}{1}% zeros-to-ramp-sums
\chapters@set{20261125T001100}{2}% inert-primes-and-estimate-b
\chapters@set{20261125T005000}{3}% galois-fourier-and-averaging
\chapters@set{20261125T005500}{4}% columns-from-a-sphere
\chapters@set{20261125T000000}{5}% slope-method
\chapters@set{20261126T030100}{6}% a-twist-for-each-prime
\chapters@set{20261125T003000}{7}% the-chain-and-the-square-root
\chapters@set{20261125T010000}{8}% rank-lemma-and-zero-estimate
\chapters@set{20261125T010500}{9}% poincare-series-and-near-orthogonality
\chapters@set{20261124T223000}{10}% elliptic-curves-for-the-proof
\chapters@set{20261125T009500}{11}% fock-model-and-theta-intertwiners
\chapters@set{20261124T223100}{12}% complex-multiplication-and-inert-reduction
\chapters@set{20261125T008000}{13}% one-cm-determinant
\chapters@set{20261124T214500}{14}% supersingular-theta-divisibility
\chapters@set{20261124T214600}{15}% the-isotropic-plane
\chapters@set{20261125T013700}{16}% pairings-on-a-formal-group-of-height-two
\chapters@set{20261125T013500}{17}% hecke-nondegeneracy-on-the-tree
\chapters@set{20261126T030200}{18}% the-bound-is-attained
\chapters@set{20261125T009000}{19}% kronecker-limit-and-chowla-selberg
\chapters@set{20261125T002000}{20}% definite-theta-lifts
\chapters@set{20261125T015500}{21}% siegel-weil-in-the-smallest-case
\chapters@set{20261125T002100}{22}% definite-theta-lifts-and-norms
\chapters@set{20261125T016500}{23}% hecke-operators-and-theta-lifts
\chapters@set{20261125T016700}{24}% divisibility-along-the-chain
\chapters@set{20261125T004000}{25}% denominators-and-sha
\chapters@set{20261125T004100}{26}% tate-twists-and-sha
\chapters@set{20261126T020430}{27}% the-minus-part-by-gauss-sums
\chapters@set{20261125T014100}{28}% congruence-modules
\chapters@set{20261125T014500}{29}% eisenstein-congruences-and-selmer-classes
\chapters@set{20261125T014700}{30}% pseudocharacters-and-the-length-bound
\chapters@set{20261125T011000}{31}% galois-representations-black-box
\chapters@set{20261126T030300}{32}% the-data
\chapters@set{20261126T020210}{33}% bessel-poincare-series-and-their-gram-operators
\chapters@set{20261126T030400}{34}% the-seesaw-identity
\chapters@set{20261126T030500}{35}% divisibility-of-the-composite-lift
\chapters@set{20261126T030600}{36}% theorem-a-short-forms
\chapters@set{20261126T030700}{37}% the-determinant-of-a-chain
\chapters@set{20261126T030800}{38}% the-cost-of-one-lift
\chapters@set{20261126T030900}{39}% denominators-of-projections-to-lifts
\chapters@set{20261126T031000}{40}% twists-with-few-congruences
\chapters@set{20261126T031100}{41}% theorem-c-and-the-main-inequality
\providecommand\chapnumber[1]{%
  \@ifundefined{chapters@#1}%
    {\PackageError{chapters}{Sheet #1 is not listed in the index}%
      {List it in 20261125T012000--grh-workbook__guide.org\MessageBreak
       and run scripts/make-chapters.py.}}%
    {\edef\thechap{\@nameuse{chapters@#1}}}}
\makeatother
\chapnumber{<sheet ID>}
\makeatletter
\providecommand\chapters@set[2]{\@namedef{chapters@#1}{#2}}
\chapters@set{20261125T001000}{1}% zeros-to-ramp-sums
\chapters@set{20261125T001100}{2}% inert-primes-and-estimate-b
\chapters@set{20261125T005000}{3}% galois-fourier-and-averaging
\chapters@set{20261125T005500}{4}% columns-from-a-sphere
\chapters@set{20261125T000000}{5}% slope-method
\chapters@set{20261126T030100}{6}% a-twist-for-each-prime
\chapters@set{20261125T003000}{7}% the-chain-and-the-square-root
\chapters@set{20261125T010000}{8}% rank-lemma-and-zero-estimate
\chapters@set{20261125T010500}{9}% poincare-series-and-near-orthogonality
\chapters@set{20261124T223000}{10}% elliptic-curves-for-the-proof
\chapters@set{20261125T009500}{11}% fock-model-and-theta-intertwiners
\chapters@set{20261124T223100}{12}% complex-multiplication-and-inert-reduction
\chapters@set{20261125T008000}{13}% one-cm-determinant
\chapters@set{20261124T214500}{14}% supersingular-theta-divisibility
\chapters@set{20261124T214600}{15}% the-isotropic-plane
\chapters@set{20261125T013700}{16}% pairings-on-a-formal-group-of-height-two
\chapters@set{20261125T013500}{17}% hecke-nondegeneracy-on-the-tree
\chapters@set{20261126T030200}{18}% the-bound-is-attained
\chapters@set{20261125T009000}{19}% kronecker-limit-and-chowla-selberg
\chapters@set{20261125T002000}{20}% definite-theta-lifts
\chapters@set{20261125T015500}{21}% siegel-weil-in-the-smallest-case
\chapters@set{20261125T002100}{22}% definite-theta-lifts-and-norms
\chapters@set{20261125T016500}{23}% hecke-operators-and-theta-lifts
\chapters@set{20261125T016700}{24}% divisibility-along-the-chain
\chapters@set{20261125T004000}{25}% denominators-and-sha
\chapters@set{20261125T004100}{26}% tate-twists-and-sha
\chapters@set{20261126T020430}{27}% the-minus-part-by-gauss-sums
\chapters@set{20261125T014100}{28}% congruence-modules
\chapters@set{20261125T014500}{29}% eisenstein-congruences-and-selmer-classes
\chapters@set{20261125T014700}{30}% pseudocharacters-and-the-length-bound
\chapters@set{20261125T011000}{31}% galois-representations-black-box
\chapters@set{20261126T030300}{32}% the-data
\chapters@set{20261126T020210}{33}% bessel-poincare-series-and-their-gram-operators
\chapters@set{20261126T030400}{34}% the-seesaw-identity
\chapters@set{20261126T030500}{35}% divisibility-of-the-composite-lift
\chapters@set{20261126T030600}{36}% theorem-a-short-forms
\chapters@set{20261126T030700}{37}% the-determinant-of-a-chain
\chapters@set{20261126T030800}{38}% the-cost-of-one-lift
\chapters@set{20261126T030900}{39}% denominators-of-projections-to-lifts
\chapters@set{20261126T031000}{40}% twists-with-few-congruences
\chapters@set{20261126T031100}{41}% theorem-c-and-the-main-inequality
\providecommand\chapnumber[1]{%
  \@ifundefined{chapters@#1}%
    {\PackageError{chapters}{Sheet #1 is not listed in the index}%
      {List it in 20261125T012000--grh-workbook__guide.org\MessageBreak
       and run scripts/make-chapters.py.}}%
    {\edef\thechap{\@nameuse{chapters@#1}}}}
\makeatother

\ifcsname chapters@20261126T070000\endcsname
  \chapnumber{20261126T070000}
\else
  \renewcommand{\thechap}{?}
\fi
\renewcommand{\thesection}{\thechap.\arabic{section}}

\theoremstyle{definition}
\newtheorem{problem}{Problem}[section]

\newcommand{\Z}{\mathbb{Z}}
\newcommand{\Q}{\mathbb{Q}}
\newcommand{\C}{\mathbb{C}}
\newcommand{\F}{\mathbb{F}}
\newcommand{\Oc}{\mathcal{O}}
\newcommand{\Cg}{\mathrm{C}}
\newcommand{\pr}{\mathrm{pr}}
\newcommand{\len}{\operatorname{len}}
\newcommand{\GL}{\mathrm{GL}}
\newcommand{\Frob}{\mathrm{Fr}}
\newcommand{\cyc}{\chi_{\mathrm{cyc}}}
\newcommand{\val}{\operatorname{val}}
\newcommand{\ord}{\operatorname{ord}}

\setlength{\emergencystretch}{3em}
\title{Clearing a projection: two classes}
\author{}
\date{November 26, 2026}

\begin{document}
\maketitle
\refstepcounter{chap}\label{chap}

\paragraph{Prerequisites.}
Discrete valuation rings and lengths of modules; the elementary divisor theorem.  The Chebotarev density theorem and the Brauer--Nesbitt theorem \cite[(30.16)]{CR}.  From Chapter~\ref{cg-chap}, \emph{Congruence modules}: the congruence module \(\Cg(\,\cdot\,;\,\cdot\,)\) and Problems~\ref{cg-pr:symmetry} and~\ref{cg-pr:denominator}.  From Chapter~\ref{ec-chap}, \emph{Eisenstein congruences and Selmer classes}: Problem~\ref{ec-pr:diag} and Ribet's lemma, Theorem~\ref{ec-thm:ribet-lemma}.  From Chapter~\ref{hk-chap}, \emph{Hecke operators and theta lifts}: Section~\ref{hk-sec:brandt}, with Problems~\ref{hk-pr:brandt} and~\ref{hk-pr:basis}.  From Chapter~\ref{gr-chap}: Section~\ref{gr-sec:es} on \(X_0(11)\).  For Section~\ref{sec:classes}: Chapters~\ref{tt-chap} and~\ref{gs-chap} on Tate twists and odd classes, the chosen places and the integer \(\mathcal N\) of Section~\ref{ob-sec:chosen}, and the statement of Chapter~\ref{dp-chap}, \emph{Denominators of projections to lifts}; Chapter~\ref{pc-chap} is background.

Chapter~\ref{dp-chap} states how large the denominators of orthogonal projections onto spaces of lifts can be, measured by the determinant of an operator that clears them.  This sheet illustrates the denominator in a model with two quaternionic ideal classes.  In Chapter~\ref{cg-chap} the exponent of a congruence module gave the scalar denominator of a projection onto an eigenspace, and in Chapter~\ref{ec-chap} a congruence produced a Galois cohomology class.  Section~\ref{sec:clearing} measures a denominator by a clearing operator and shows why false lifts must be excluded; Section~\ref{sec:eleven} computes the denominator for a class set with two elements, where it is \(5\); Section~\ref{sec:classes} explains the congruence behind the denominator and the costs on the Galois side that enter the clearing bound.  In Sections~\ref{sec:clearing} and~\ref{sec:eleven} the letters \(K\), \(T\), \(V\), \(L\), \(Y\) and \(\Pi\) are local.

\section{Clearing a projection}\label{sec:clearing}

Let \(\Oc\) be a discrete valuation ring with fraction field \(K\) and valuation \(\ord\), normalized by \(\ord(\varpi) = 1\) for a uniformizer \(\varpi\); we use \(\ord(0) = +\infty\).  Let \(V\) be a finite-dimensional \(K\)-vector space with a lattice \(L\), and let \(V = \mathcal F \oplus \mathcal F'\) with projection \(\pr \colon V \to \mathcal F\) along \(\mathcal F'\).  Write \(\Cg(L; \mathcal F) := \pr(L)/(L \cap \mathcal F)\), as in Chapter~\ref{cg-chap}.  An operator \(C_\ell \in \mathrm{End}_K(\mathcal F)\) (the subscript recalls the prime \(\ell\) of the later sections) \emph{clears} \(\pr\) on a subset \(Y \subseteq L\) if \(C_\ell(\pr(y)) \in L\) for every \(y \in Y\).  A scalar \(c\) clears \(\pr\) on \(L\) when \(c\,\pr(L) \subseteq L\); an operator may clear different directions of \(\mathcal F\) by different amounts, and we measure it by \(\ord(\det C_\ell)\).

\begin{problem}\label{pr:clear-det}
\(C_\ell\) clears \(\pr\) on \(L\) if and only if \(C_\ell(\pr(L)) \subseteq L \cap \mathcal F\).  The least value of \(\ord(\det C_\ell)\) over such \(C_\ell\) is \(\len_{\Oc} \Cg(L; \mathcal F)\).
\end{problem}

An integral operator that is a nonzero scalar \(a\) on \(\mathcal F\) and \(0\) on \(\mathcal F'\) clears \(\pr\) on \(L\) as the scalar \(a\), so \(\len_{\Oc}\Cg(L; \mathcal F) \le \dim \mathcal F \cdot \ord(a)\); this is \cite[Lemma~5.1]{Zh}.  In part (b) of the next problem, an eigenvalue difference gives the exact denominator.

\begin{problem}\label{pr:proj-cost}
(a) The orthogonal projection of \(\Q^2\) onto the line through \((1,2)\) is \(\frac15\left(\begin{smallmatrix}1&2\\2&4\end{smallmatrix}\right)\), and, for \(c \in \Q\), \(c\) times it preserves \(\Z^2\) exactly when \(c \in 5\Z\).  (b) Let \(s, c \in \Z\) with \(s \ge 1\), and \(T = \left(\begin{smallmatrix}\ell^s&1\\0&0\end{smallmatrix}\right)\) over \(\Z_\ell\).  The projection onto the \(\ell^s\)-eigenline along \(\ker T\) is \(\ell^{-s}T\), and \(\ell^{c-s}T\) preserves \(\Z_\ell^2\) exactly when \(c \ge s = \val_\ell(\ell^s - 0)\).
\end{problem}

The operators that clear \(\pr\) on a given \(Y\) form an \(\Oc\)-module; the definitions make sense for any valuation ring \(\Oc\) and any \(\Oc\)-submodule \(L\) of \(V\).  The next problem combines clearing operators that are diagonal in one basis.

\begin{problem}\label{pr:diagonal}
Let \(\Oc\) be a valuation ring with infinite residue field and valuation \(\ord\), such as the ring of integers of \(\bar{\Q}_\ell\) with \(\ord(\ell) = 1\), and let \(b_1, \dots, b_n\) be a basis of \(\mathcal F\).  Let \(C_\ell^{(1)}, \dots, C_\ell^{(n)}\) clear \(\pr\) on \(Y\), with \(C_\ell^{(j)} b_i = \lambda_{ji} b_i\).  Then some \(C_\ell = \sum_j u_j C_\ell^{(j)}\) with \(u_j \in \Oc\) has \(C_\ell b_i = \mu_i b_i\) with \(\ord(\mu_i) \le \ord(\lambda_{ii})\) for every \(i\); it clears \(\pr\) on \(Y\).
\end{problem}

Now let \(K\) be a subfield of \(\C\) stable under complex conjugation, and \(\Oc\) a discrete valuation ring with fraction field \(K\).  Let \(V\) carry a positive definite hermitian form \(\langle\,,\rangle\), linear in the first variable, let \(\mathcal F' = \mathcal F^\perp\), and let \(T\) be a self-adjoint operator on \(V\) with \(T(L) \subseteq L\), diagonalizable with eigenvalues in \(\Oc\), such that \(\mathcal F \subseteq \ker(T - \lambda)\).  Call the vectors of \(\mathcal F^\perp\) with the eigenvalue of \(\mathcal F\),
\[
  F_T := \mathcal F^\perp \cap \ker(T - \lambda) ,
\]
the \emph{false lifts} for \(T\).  Let \(\Pi(X) := \prod_\mu (X - \mu)\), the product over the eigenvalues \(\mu \ne \lambda\) of \(T\), each taken once.

\begin{problem}\label{pr:hecke-clear}
For \(y \in L\) orthogonal to \(F_T\), \(\Pi(T)\,y = \Pi(\lambda)\,\pr(y)\).  Hence the scalar \(\Pi(\lambda)\) clears \(\pr\) on the \(y \in L\) orthogonal to \(F_T\), and the determinant of the scalar \(\Pi(\lambda)\) acting on \(\mathcal F\) has valuation \(\dim \mathcal F \cdot \ord(\Pi(\lambda))\).
\end{problem}

The cost \(\ord(\Pi(\lambda)) = \sum_\mu \ord(\lambda - \mu)\) counts congruences between the eigenvalue of \(\mathcal F\) and the other eigenvalues.  The restriction to vectors orthogonal to the false lifts cannot be dropped.

\begin{problem}\label{pr:false-needed}
Let \(V = K^2\) with \(\langle x, y \rangle = x_1\bar y_1 + x_2\bar y_2\), \(L = \Oc^2\), \(T = \lambda\cdot 1\), \(\mathcal F = K(1,1)\), and suppose \(\ord(2) > 0\).  Then \(F_T = K(1,-1)\), no unit \(c \in \Oc^\times\) clears \(\pr\) on \(L\), and \(c = 1\) clears it on the \(y \in L\) orthogonal to \(F_T\).
\end{problem}

\section{Two classes}\label{sec:eleven}

We take the class set of Section~\ref{hk-sec:brandt} of Chapter~\ref{hk-chap}: the two left ideal classes of a maximal order in the quaternion algebra over \(\Q\) ramified at \(11\) and \(\infty\), indexed by their right orders, with \(w_1 = 2\) and \(w_2 = 3\).  Functions on the two classes are vectors \(\varphi = (\varphi(1), \varphi(2))\), with
\[
  \langle \varphi, \varphi' \rangle := \frac{\varphi(1)\overline{\varphi'(1)}}{w_1} + \frac{\varphi(2)\overline{\varphi'(2)}}{w_2} ,
\]
and the Brandt matrix \(B(\ell)\), for a prime \(\ell \ne 11\), acts by \((B(\ell)\varphi)(i) = \sum_j B_{ij}(\ell)\varphi(j)\).  Recall from that section that \(B(\ell)\) has integer entries, that \(w_jB_{ij}(\ell) = w_iB_{ji}(\ell)\) and \(\sum_j B_{ij}(\ell) = \ell + 1\), and that \(B(\ell)\) is the action of \(T_\ell\) on \(M_2(\Gamma_0(11))\), which is spanned by an Eisenstein series of \(T_\ell\)-eigenvalue \(\ell + 1\) and by the form called \(f\) there; we write
\[
  f_{11} := q\prod_{n \ge 1}(1-q^n)^2(1-q^{11n})^2 = \sum_{n \ge 1} a_n q^n .
\]
By Problem~\ref{hk-pr:brandt}, \(B(2) = \begin{pmatrix} 1 & 2 \\ 3 & 0 \end{pmatrix}\) and \(B(3) = \begin{pmatrix} 2 & 2 \\ 3 & 1 \end{pmatrix}\).  The constant function \(1\) plays the role of the lifts: \(\mathcal F\) is the line of constant functions, and \(\pr\) is the orthogonal projection onto \(\mathcal F\).  The integral functions are \(L := \Z^2\).  To apply Section~\ref{sec:clearing} over \(\Z_{(\ell)}\), use the lattice \(\Z_{(\ell)}^2\) in \(\Q^2\).  Over \(\Q_\ell\) use the lattice \(\Z_\ell^2\) and the same bilinear formula without the bar.

\begin{problem}\label{pr:eleven-adjoint}
\(B(\ell)\) is self-adjoint for \(\langle\,,\rangle\), \(B(\ell)1 = (\ell + 1)1\), and \(\varphi_{11} := (2, -3)\) spans \(\mathcal F^\perp\) and satisfies \(B(\ell)\varphi_{11} = a_\ell\varphi_{11}\).
\end{problem}

\begin{problem}\label{pr:eleven-pr}
\(\pr(\varphi) = \frac15\bigl(3\varphi(1) + 2\varphi(2)\bigr)\cdot 1\), and \(5\,\pr = B(2) + 2 = B(3) + 1\).  Hence \(\pr(L) = \frac15\Z\cdot 1\) and \(\Cg(L; \mathcal F) \cong \Z/5\Z\).  The least value of \(\val_5(\det C_\ell)\) for \(C_\ell\) clearing \(\pr\) on \(\Z_5^2\) is \(1\); for \(\ell \ne 5\) the identity clears \(\pr\) on \(\Z_\ell^2\).  The scalar \(\Pi(\lambda)\) of Problem~\ref{pr:hecke-clear} is \(5\) both for \(T = B(2)\) and for \(T = B(3)\).
\end{problem}

What congruence do the Brandt eigenvalues satisfy?

\begin{problem}\label{pr:eleven-cong}
\(B(\ell) - a_\ell = (\ell + 1 - a_\ell)\pr\).  Hence \(a_\ell \equiv 1 + \ell \pmod 5\) for every prime \(\ell \ne 11\); for example \(a_{13} = 4\) (Problem~\ref{gr-pr:coeffs}).
\end{problem}

The denominator \(5\) is the numerator of the mass \(1/w_1 + 1/w_2 = (11 - 1)/12 = 5/6\).  The next problem gives the denominator at primes dividing none of the weights.

\begin{problem}\label{pr:mass}
Let \(I\) be a nonempty finite set with positive integers \(w_i\), let \(\langle \varphi, \varphi' \rangle := \sum_i \varphi(i)\varphi'(i)/w_i\) on \(\Q_\ell^I\), let \(\mathcal F\) be the line of constant functions, and let \(\ell\) be a prime dividing no \(w_i\).  Then \(\mathfrak m := \sum_i 1/w_i\) has \(\val_\ell(\mathfrak m) \ge 0\), and \(\Cg(\Z_\ell^I; \mathcal F) \cong \Z/\ell^{\val_\ell(\mathfrak m)}\Z\).
\end{problem}

For the quaternion algebra ramified at a prime \(p\) and \(\infty\) the mass is \((p-1)/12\), and Mazur showed that the index of his Eisenstein ideal is the numerator of \((p-1)/12\) \cite{Ma}.  There is also an integrality question for the pairing of a cleared form with an integral form.

\begin{problem}\label{pr:pairing}
In the setting of Problem~\ref{pr:mass}, drop the condition on \(\ell\), and let \(n_0\) be a common multiple of the \(w_i\).  If \(C_\ell\) clears \(\pr\) on \(Y \subseteq \Z_\ell^I\), then so does \(n_0C_\ell\), and \(\langle n_0C_\ell\pr(y), \varphi' \rangle \in \Z_\ell\) for \(y \in Y\) and \(\varphi' \in \Z_\ell^I\).  In the class set with two elements, \(6\langle \varphi, \varphi' \rangle \in \Z\) for \(\varphi, \varphi' \in \Z^2\).
\end{problem}

\section{Lifts and the classes they produce}\label{sec:classes}

Where does the congruence of Problem~\ref{pr:eleven-cong} come from?  The constant function carries the Galois representation \(1 \oplus \cyc\), where \(\cyc\) is the \(5\)-adic cyclotomic character; its trace at an arithmetic Frobenius \(\Frob_\ell\) is \(1 + \ell\) for \(\ell \ne 5\).  The form \(f_{11}\) carries the representation \(\rho_{11}\) of \(\mathrm{Gal}(\bar{\Q}/\Q)\) on the rational \(5\)-adic Tate module \(T_5(X_0(11)) \otimes \Q_5\), whose trace at \(\Frob_\ell\) is \(a_\ell\) for \(\ell \nmid 55\) (Section~\ref{gr-sec:es} of Chapter~\ref{gr-chap}).  Let \(\omega\) be the reduction of \(\cyc\) modulo \(5\).

\begin{problem}\label{pr:residual}
The semisimplification of the reduction of \(\rho_{11}\) modulo \(5\) is \(1 \oplus \omega\), and \(\omega(c) = -1\) for every complex conjugation \(c\).
\end{problem}

The representation \(\rho_{11}\) is irreducible: it is semisimple and its commutant is \(\Q_5\), by Faltings' theorem \cite{Fa}, since \(X_0(11)\) has no complex multiplication (its \(j\)-invariant \(-122023936/161051\) is not an integer).  Ribet's lemma, Theorem~\ref{ec-thm:ribet-lemma}, then gives a nonzero class with values in \(\F_5(\omega)\) or \(\F_5(\omega^{-1})\), an odd character.

In Chapter~\ref{dp-chap}, \(E\) is a CM field with maximal real subfield \(E^+\).  Suppose \(R\) is a main summand for a descent pair and an eigenvalue system in the sense of Definition~\ref{dp-def:main-summand}.  For every compatible pair twisted by a tame character \(\xi\), every eigenform lifting that system, and every representation attached to its Hecke polynomials, the trace of the normalized representation at \(g\in\mathrm{Gal}(\bar{\Q}/E)\) is \(\hat\xi(g)\operatorname{tr}R(g)+\sum_{0\le e<t'}\cyc(g)^{-e}\).  Here \(\hat\xi\) is the avatar of \(\xi^{-1}\), \(t'\) is the difference between the target and source ranks, and \(\cyc\) is the \(\ell\)-adic cyclotomic character.  For an extension of \(\cyc^{-e'}\) by \(\cyc^{-e}\) modulo \(\ell^\nu\), the extension cocycle takes values in \((\Oc/\ell^\nu)(-a)\), where \(a = e - e'\).  On cocycles \(z\) of \(\mathrm{Gal}(\bar{\Q}/E)\) a complex conjugation \(c\) acts by \(z^c(g) := z(cgc)\); call \(z\) \emph{odd} if \(z^c = -z\) and \emph{even} if \(z^c = z\).

\begin{problem}\label{pr:parity}
Let \(\ell\) be odd and \(\Oc\) the ring of integers of a finite extension of \(\Q_\ell\).  Every cocycle of \(\mathrm{Gal}(\bar{\Q}/E^+)\) with values in \((\Oc/\ell^\nu)(-a)\) differs by a coboundary from one with \(z(c) = 0\).  For such \(z\), the restriction to \(\mathrm{Gal}(\bar{\Q}/E)\) satisfies \(z^c = (-1)^a z\).
\end{problem}

For odd \(\ell\) and odd \(a\), restriction from \(\mathrm{Gal}(\bar{\Q}/E^+)\) to \(\mathrm{Gal}(\bar{\Q}/E)\) therefore gives classes with odd representatives.  Apply Theorem~\ref{tt-thm:oddsha} with \(N=q_1d\).  It gives one positive integer \(X_{\mathrm{odd}}\) for all primes \(\ell\), all coefficient rings \(\Oc\) specified in that theorem, all \(0\le a<t\) and all \(\nu\ge0\).  For \(a>0\), \(\val_{\Oc}(X_{\mathrm{odd}})\) bounds the \(\Oc\)-length of the module of odd cocycles with values in \((\Oc/\ell^\nu)(-a)\) and the stated local conditions.  For \(a=0\) it bounds the \(\Oc\)-length of the module of all continuous everywhere-unramified homomorphisms from \(\mathrm{Gal}(\bar{\Q}/E)\) to \(\Oc/\ell^\nu\).  The proof for positive \(a\) uses Gauss sums as in Chapter~\ref{gs-chap}.  For \(0<a<t'\) the proof controls the even cocycles by tests at auxiliary primes (Chapter~\ref{n2-chap}).  The cost of these tests enters the clearing bound of Chapter~\ref{dp-chap}, which contains no separate term for the length of a module of even cocycles.  A class can be tested by its values at Frobenius elements; here is an example, for a Kummer class.

\begin{problem}\label{pr:kummer}
Let \(\ell\) be an odd prime, \(u \in \Q^\times\) not an \(\ell\)-th power, and \(\mathrm{kum}_u\) the Kummer class of \(u\) in \(H^1(\mathrm{Gal}(\bar{\Q}/\Q), \mu_\ell)\).  For a prime \(q \equiv 1 \pmod \ell\) with \(\val_q(u) = 0\), \(\mathrm{kum}_u\) is unramified at \(q\), and \(\mathrm{kum}_u(\Frob_q) = 0\), a condition independent of the choice of \(\Frob_q\), if and only if \(u^{(q-1)/\ell} \equiv 1 \pmod q\).  There are infinitely many such \(q\) with \(\mathrm{kum}_u(\Frob_q) \ne 0\).  For \(\ell = 3\) and \(u = 2\), \(q = 7\) detects \(\mathrm{kum}_2\) and \(q = 31\) does not.
\end{problem}

Ribet's lemma needs an element on which the two characters differ (Problem~\ref{ec-pr:diag}).  For the characters \(\cyc^{-e}\), \(0 \le e < t\), with \(\cyc\) the \(\ell\)-adic cyclotomic character, the cost of an element \(g\) is the valuations of the differences \(\cyc(g)^{a} - 1\), \(0 < a < t\).  If \(\cyc(g)\) reduces to a generator of \(\F_\ell^\times\) and \(t \le \ell - 1\), these valuations are \(0\), by Problem~\ref{pr:generator}.  If \(g\) is an arithmetic Frobenius element at a chosen place \(w\) not dividing \(\ell\) (Section~\ref{ob-sec:chosen}), then \(\cyc(g) = \#k_w\) and these valuations are at most \(\val_\ell(\mathcal N)\), by Problem~\ref{pr:gl-order} with \(n = rt\): the integer \(\mathcal N\) of that section is a multiple of \(\#\GL_{rt}(k_w)\).

\begin{problem}\label{pr:generator}
Let \(g \in \Z_\ell^\times\) reduce to a generator of \(\F_\ell^\times\).  Then \(g^a - 1 \in \Z_\ell^\times\) for \(0 < a < \ell - 1\).
\end{problem}

\begin{problem}\label{pr:gl-order}
Let \(q\) be a prime power and \(1 \le a \le n\).  Then \(q^a - 1\) divides \(\#\GL_n(\F_q) = q^{n(n-1)/2}\prod_{i=1}^n (q^i - 1)\).
\end{problem}


\paragraph{In the Lean formalization.}
In the Lean proof \cite{GL} (\nolinkurl{grh/src/}, with the namespace \texttt{New.} omitted), the trace identity of Section~\ref{sec:classes} occurs in \texttt{Selmer.IsMainSummand} and \texttt{Selmer.HasLiftTrace}; the uniform length bound is \texttt{exists\_length\_shaOdd\_le}, whose witness for \(N=q_1d\) is \texttt{Selmer.oddClassBound}; the auxiliary primes come from \texttt{Selmer.exists\_isPrimeTriangle}, and \texttt{Selmer.hasTriangleCost\_tateTwist} applies their tests to the even cocycles; the separation bound is \texttt{Selmer.exists\_cyclotomicChar\_pow\_sub\_one\_dvd}, with \(\mathcal N\) the integer \texttt{Data.chosenGroupOrder}.

\paragraph{What this sheet takes as given.}
Gross's description of the Brandt matrices at \(11\) (Section~\ref{hk-sec:brandt}); the Tate module of \(X_0(11)\) and its Frobenius traces (Section~\ref{gr-sec:es}); Faltings' theorem; Ribet's lemma (Theorem~\ref{ec-thm:ribet-lemma}); the Brauer--Nesbitt and Chebotarev theorems; the odd-class bound, Theorem~\ref{tt-thm:oddsha}; and Mazur's theorem, quoted only as a remark.

\section*{Hints}

\subsection*{\S\ref{sec:clearing}}
\begin{description}
\item[\ref{pr:clear-det}] \(C_\ell\) maps \(\mathcal F\) to \(\mathcal F\).  For the minimum, \(\pr(L)\) and \(L \cap \mathcal F\) are lattices in \(\mathcal F\); use elementary divisors.
\item[\ref{pr:proj-cost}] Apply the projection to \(e_1\) in (a) and to \(e_2\) in (b).
\item[\ref{pr:diagonal}] Choose the \(u_j\) generically modulo the maximal ideal.
\item[\ref{pr:hecke-clear}] Decompose \(y\) into eigenvectors of \(T\); \(\ker(T - \lambda) = \mathcal F \oplus F_T\) is orthogonal to the other eigenspaces.
\end{description}

\subsection*{\S\ref{sec:eleven}}
\begin{description}
\item[\ref{pr:eleven-adjoint}] Apply the balance and row-sum identities to the two orthogonal lines.
\item[\ref{pr:eleven-pr}] Apply the projection onto constants to the coordinate functions; then Problem~\ref{pr:clear-det}.
\item[\ref{pr:eleven-cong}] Compare both sides on a basis.
\item[\ref{pr:mass}] Compute \(\pr\) on the indicator functions.
\item[\ref{pr:pairing}] Inspect each weight in the class sum.
\end{description}

\subsection*{\S\ref{sec:classes}}
\begin{description}
\item[\ref{pr:residual}] The characteristic polynomial of \(\rho_{11}(\Frob_\ell)\) is \(X^2 - a_\ell X + \ell\) for \(\ell \nmid 55\) (Section~\ref{gr-sec:es} of Chapter~\ref{gr-chap}).  Reduce on the lattice \(T_5(X_0(11))\); \(\rho_{11}\) is unramified outside \(55\).  Use Problem~\ref{pr:eleven-cong}, the Chebotarev density theorem and the Brauer--Nesbitt theorem for characteristic polynomials.
\item[\ref{pr:parity}] Evaluate the cocycle relation on \(c \cdot c = 1\) and on \(c \cdot g \cdot c\), using \(\cyc(c) = -1\).
\item[\ref{pr:kummer}] Consider the action on \(u^{1/\ell}\) at a prime above \(q\).  For infinitude, compare \(\Q(u^{1/\ell})\) with \(\Q(\mu_\ell)\) and use the Chebotarev density theorem in their compositum.
\end{description}

\begin{thebibliography}{GL}
\bibitem[CR]{CR} C.~W.~Curtis and I.~Reiner, \emph{Representation Theory of Finite Groups and Associative Algebras}, Interscience, 1962.
\bibitem[Fa]{Fa} G.~Faltings, \emph{Endlichkeitss\"atze f\"ur abelsche Variet\"aten \"uber Zahlk\"orpern}, Invent.\ Math.\ \textbf{73} (1983), 349--366.
\bibitem[GL]{GL} L.~Sernau, \emph{GRH for Dirichlet \(L\)-functions, in Lean~4}, a shorter rewriting of the Lean proof, text by Claude, \url{https://github.com/lsernau/grh-lean}, commit \texttt{72fead1}, readable edition \texttt{grh/}.
\bibitem[Ma]{Ma} B.~Mazur, \emph{Modular curves and the Eisenstein ideal}, Publ.\ Math.\ IH\'ES \textbf{47} (1977), 33--186.
\bibitem[Zh]{Zh} R.~Zhang, \emph{A preliminary note on Alp\"oge's proposed proof of the generalized Riemann hypothesis}, Hexagon v1, 21 November 2026, \url{https://hexagonmath.org/pdf/2611.00809v1}.
\end{thebibliography}

\end{document}
